\documentclass[12pt,a4paper]{article}

\usepackage[utf8]{inputenc}
\usepackage[T1]{fontenc}
\usepackage[brazil]{babel}
\usepackage[a4paper,margin=2.cm]{geometry}
\usepackage{enumitem}
\setlist[enumerate]{itemsep=0.45em,topsep=0.35em}

\begin{document}

\begin{center}
  {\Large\bfseries PROVA SIMULADA -- NPOR -- VERSÃO 3}\par
  \vspace{0.4cm}
  Nome: \rule{9cm}{0.4pt}\par\vspace{0.3cm}
  Data: \rule{4cm}{0.4pt}\hfill Nota: \rule{3cm}{0.4pt}
\end{center}

\noindent\textbf{Duração sugerida: 90 minutos.} A prova oficial pode ter outro tempo e outros conteúdos; confirme as instruções da sua Organização Militar. Este simulado tem 13 questões objetivas e uma redação.

\medskip
\noindent\textbf{Instruções:} marque uma alternativa por questão. Na redação, escreva um texto dissertativo-argumentativo, com letra legível, respeitando o tema e o espaço disponível.

\subsection*{Português}
Leia o texto para responder às questões 1 e 2.

\begin{quote}
As ferramentas digitais facilitam o acesso a informações e permitem que as pessoas aprendam em diferentes lugares. Embora sejam úteis, elas não garantem que todo conteúdo encontrado seja correto. Comparar fontes e verificar a data de publicação são atitudes importantes para formar uma opinião responsável.
\end{quote}

\begin{enumerate}
  \item Qual é a principal recomendação apresentada no texto?
  \begin{enumerate}[label=\alph*)]
    \item Evitar completamente o uso de ferramentas digitais.
    \item Aceitar como verdadeiro todo conteúdo publicado na internet.
    \item Comparar fontes e verificar informações antes de formar uma opinião.
    \item Consultar somente conteúdos sem data de publicação.
  \end{enumerate}

  \item No texto, a palavra ``Embora'' introduz uma ideia de:
  \begin{enumerate}[label=\alph*)]
    \item causa.
    \item concessão.
    \item conclusão.
    \item finalidade.
  \end{enumerate}

  \item Assinale a frase com concordância verbal adequada.
  \begin{enumerate}[label=\alph*)]
    \item Faltam duas semanas para a avaliação.
    \item Houveram várias notícias sobre o assunto.
    \item Existe bons motivos para conferir a fonte.
    \item As informações foi divulgada ontem.
  \end{enumerate}

  \item Assinale a frase correta quanto à regência verbal.
  \begin{enumerate}[label=\alph*)]
    \item Os estudantes assistiram ao debate.
    \item Os estudantes assistiram o debate, no sentido de acompanhar.
    \item A equipe obedeceu o regulamento.
    \item O candidato prefere mais estudar do que descansar.
  \end{enumerate}
\end{enumerate}

\subsection*{Matemática e raciocínio lógico}
\begin{enumerate}[resume]
  \item O preço de um produto passou de R\$ 240,00 para R\$ 300,00. Qual foi o aumento percentual?
  \begin{enumerate}[label=\alph*)]
    \item 20\%
    \item 25\%
    \item 30\%
    \item 60\%
  \end{enumerate}

  \item Cinco litros de tinta cobrem 40 m$^2$. Mantendo o mesmo rendimento, quantos metros quadrados podem ser cobertos com 7,5 litros?
  \begin{enumerate}[label=\alph*)]
    \item 50 m$^2$
    \item 55 m$^2$
    \item 60 m$^2$
    \item 75 m$^2$
  \end{enumerate}

  \item A equação $x^2-5x+6=0$ tem duas raízes. Qual é a menor delas?
  \begin{enumerate}[label=\alph*)]
    \item 1
    \item 2
    \item 3
    \item 6
  \end{enumerate}

  \item Um triângulo tem base de 12 cm e altura de 5 cm. Qual é sua área?
  \begin{enumerate}[label=\alph*)]
    \item 17 cm$^2$
    \item 30 cm$^2$
    \item 60 cm$^2$
    \item 120 cm$^2$
  \end{enumerate}

  \item Complete a sequência: 2, 5, 10, 17, \ldots
  \begin{enumerate}[label=\alph*)]
    \item 24
    \item 25
    \item 26
    \item 27
  \end{enumerate}
\end{enumerate}

\subsection*{Conhecimentos gerais e institucionais}
\begin{enumerate}[resume]
  \item Em que ano Brasília foi inaugurada como capital do Brasil?
  \begin{enumerate}[label=\alph*)]
    \item 1945
    \item 1950
    \item 1960
    \item 1988
  \end{enumerate}

  \item Qual região brasileira possui o maior número de estados?
  \begin{enumerate}[label=\alph*)]
    \item Norte.
    \item Nordeste.
    \item Centro-Oeste.
    \item Sul.
  \end{enumerate}

  \item Qual expressão aparece na Bandeira Nacional brasileira?
  \begin{enumerate}[label=\alph*)]
    \item União e Justiça.
    \item Ordem e Progresso.
    \item Pátria e Liberdade.
    \item Paz e Igualdade.
  \end{enumerate}

  \item Qual posto vem imediatamente depois de Primeiro-Tenente na hierarquia do oficialato?
  \begin{enumerate}[label=\alph*)]
    \item Aspirante-a-Oficial.
    \item Capitão.
    \item Major.
    \item Segundo-Tenente.
  \end{enumerate}
\end{enumerate}

\newpage
\section*{Redação}
\noindent\textbf{Tema:} O papel dos jovens no combate à desinformação e na construção de uma sociedade mais consciente.

\medskip
\noindent Escreva um texto dissertativo-argumentativo, apresentando seu ponto de vista, argumentos que o sustentem e uma conclusão. Sugestão: de 15 a 25 linhas. Não é necessário escrever um título.

\vspace{0.4cm}
\noindent\rule{\textwidth}{0.3pt}\par\vspace{0.45cm}
\noindent\rule{\textwidth}{0.3pt}\par\vspace{0.5cm}
\noindent\rule{\textwidth}{0.3pt}\par\vspace{0.5cm}
\noindent\rule{\textwidth}{0.3pt}\par\vspace{0.5cm}
\noindent\rule{\textwidth}{0.3pt}\par\vspace{0.5cm}
\noindent\rule{\textwidth}{0.3pt}\par\vspace{0.5cm}
\noindent\rule{\textwidth}{0.3pt}\par\vspace{0.5cm}
\noindent\rule{\textwidth}{0.3pt}\par\vspace{0.5cm}
\noindent\rule{\textwidth}{0.3pt}\par\vspace{0.5cm}
\noindent\rule{\textwidth}{0.3pt}\par\vspace{0.5cm}
\noindent\rule{\textwidth}{0.3pt}\par\vspace{0.5cm}
\noindent\rule{\textwidth}{0.3pt}\par\vspace{0.5cm}
\noindent\rule{\textwidth}{0.3pt}\par\vspace{0.5cm}
\noindent\rule{\textwidth}{0.3pt}\par\vspace{0.5cm}
\noindent\rule{\textwidth}{0.3pt}\par\vspace{0.5cm}
\noindent\rule{\textwidth}{0.3pt}\par\vspace{0.5cm}
\noindent\rule{\textwidth}{0.3pt}\par\vspace{0.5cm}
\noindent\rule{\textwidth}{0.3pt}\par\vspace{0.5cm}
\noindent\rule{\textwidth}{0.3pt}\par\vspace{0.5cm}
\noindent\rule{\textwidth}{0.3pt}\par\vspace{0.5cm}
\noindent\rule{\textwidth}{0.3pt}\par\vspace{0.5cm}
\noindent\rule{\textwidth}{0.3pt}\par

\end{document}